\section{Resonancia de Plasmones Superficiales}

\subsection{Fenomenología Física}

La resonancia de plasmones superficiales localizados constituye un fenómeno de
acoplamiento electrodinámico. Físicamente, ocurre cuando la radiación
electromagnética incidente polariza la nanopartícula metálica, lo que genera un
dipolo oscilante. Esta interacción induce una oscilación colectiva y cuantizada
de la nube de electrones de conducción, la cual se encuentra espacialmente
confinada a la superficie de la estructura nanométrica. La frecuencia propia de
esta resonancia depende intrínsecamente de la morfología, la composición y el
tamaño de la nanopartícula, confiriéndole propiedades ópticas especificas.
Fenomenológicamente, el campo electromagnético actúa como la fuerza
impulsora sobre el plasma de electrones libres, estableciendo una condición de
resonancia fuertemente acoplada a las propiedades dieléctricas del entorno.

\subsection{Cambio de absorbencia}

En un estado de dispersión coloidal estable, las nanopartículas de oro exhiben
un pico de absorbancia característico en la vecindad de \qty{520}{\nano\metre}.
El corrimiento de este pico hacia longitudes de onda mayores, fenómeno conocido
como efecto batocrómico, surge como una consecuencia electrodinámica directa de
la agregación espacial de las nanopartículas. Al reducirse la
distancia inter-partícula debido a la formación del complejo de coordinación,
los campos locales de los dipolos plasmónicos individuales interactúan
fuertemente. Este acoplamiento electromagnético altera
el entorno dieléctrico local y disminuye la energía de
resonancia del sistema acoplado. De esta manera, la modificación de
las condiciones de frontera dieléctricas desplaza las frecuencias de absorción
hacia el extremo rojo del espectro visible, evidenciando macroscópicamente la
agregación inducida.

\subsection{Información Física}

La caracterización espectral del corrimiento plasmónico permite extraer
información física precisa sobre la dinámica del sistema coloidal. En primer
lugar, la magnitud del desplazamiento batocrómico y la aparición de un efecto
hipocrómico se correlacionan directamente con el grado de agregación física de
las nanopartículas. Al analizar la dependencia empírica entre
la longitud de onda del pico de resonancia y la concentración del agente
iónico, es posible establecer una función de transferencia empírica. Esta
función paramétrica permite cuantificar la concentración exacta del ion
analizado en la solución acuosa. Adicionalmente, el perfil
espectral y la posición inicial del pico de absorbancia proporcionan
información geométrica, permitiendo verificar el diámetro de las
nanopartículas empleadas en la sonda.
